%==========================================================
\section{Protection de pile et synchronisation}
%==========================================================

% --- Stack guard ---
\begin{frame}{Protection de pile (\texttt{USE\_STACK\_PROT})}
\begin{block}{Problématique}
        Les piles des threads sont allouées via \texttt{malloc}. Un débordement (\textit{overflow}) corrompt silencieusement la mémoire des threads voisins.
    \end{block}
\centering
\begin{tikzpicture}[
  blk/.style={draw, minimum height=1.0cm, font=\footnotesize, align=center},
  arr/.style={->, >=Stealth, thick},
]
  % Memory layout
  \node[blk, fill=blue!20,  minimum width=2.5cm] (st)  at (0,   0) {\texttt{thread\_s}\\struct};
  \node[blk, fill=red!30,   minimum width=1.5cm] (gp)  at (2.0, 0) {\texttt{GUARD}\\PAGE};
  \node[blk, fill=green!20, minimum width=4.0cm] (sk)  at (5.0, 0) {Stack\\$\longleftarrow$ croît vers le bas};

  \node[font=\tiny, gray] at (0, -0.75)   {\texttt{posix\_memalign}};
  \node[font=\tiny, red!80, align=center] at (2.0,-0.75) {\texttt{PROT\_NONE}\\(avec \texttt{mprotect()})};

  % Overflow arrow
  \draw[arr, red!80, very thick] (3.0, 0.5) -- (2.0, 0.5)
    node[midway, above, font=\scriptsize, red!80] {stack overflow};

  % SIGSEGV
  \draw[arr, red!80, dashed] (2.0, 1.0) -- (2.0, 2.0)
    node[above, font=\small, red!80] {\texttt{SIGSEGV}};

  % Handler
  \node[draw, fill=orange!20, rounded corners=3pt,
        minimum width=3.5cm, minimum height=0.7cm, font=\footnotesize]
    (h) at (6.5, 2.0) {\texttt{segfault\_handler} : message + exit};

  \draw[arr, orange!70] (2.0, 2.35) -- (h.west);

  % sigaltstack
  \node[draw, fill=teal!15, rounded corners=3pt,
        minimum width=3.8cm, minimum height=0.6cm, font=\scriptsize]
    at (6.5, -1.3)
    {\texttt{sigaltstack} : pile de secours pour le handler};

  \node[font=\scriptsize, gray, align=center] at (2.0, -1.8)
    {Où est éxecuté le handler ?};
\end{tikzpicture}


\alert{Comportement validé par :} Test 90-stack\_overflow

\end{frame}

% --- Mutex ---
\begin{frame}{Sémaphores et Mutex : Une Logique Commune}
    

    \begin{itemize}
        \item \textbf{Gestion des files d'attente} :
            \begin{itemize}
                \item Utilisation d'une \texttt{waitingQueue} de type \texttt{TAILQ}.
                \item Transition d'état identique : \texttt{RUNNING} $\rightarrow$ \texttt{BLOCKED} lors de l'attente.
            \end{itemize}
        \item \textbf{Attente Passive et Atomicité} :
            \begin{itemize}
                \item Utilisation des primitives \texttt{\_\_sync} pour garantir l'exclusion mutuelle
                \item Aucun \texttt{while} d'attente active
            \end{itemize}
        \item \textbf{Le passage de verrou} :
            \begin{itemize}
                \item \texttt{unlock} et \texttt{post} partagent la même logique : le verrou est directement transféré au prochain thread en attente sans être relâché
            \end{itemize}
    \end{itemize}
\end{frame}

\begin{frame}{Variables de Condition}
               
\begin{itemize}
    \item \textbf{Synchronisation logique} : Permet d'attendre la vérification d'une condition sans gaspiller de cycles CPU (attente passive).
    \item \textbf{Mécanique interne} :
        \begin{itemize}
            \item Libération atomique du mutex associé lors du \texttt{wait}.
            \item Mise en sommeil du thread dans une file d'attente spécifique (\texttt{waitingQueue}).
            \item Réacquisition obligatoire du mutex au réveil avant de poursuivre.
        \end{itemize}
    \item \textbf{Signal vs Broadcast} :
        \begin{itemize}
            \item \texttt{Signal} : Réveil ciblé du premier thread de la file.
            \item \texttt{Broadcast} : Libération de tous les threads pour les changements d'états globaux $\rightarrow$ Utile pour les barrières.
        \end{itemize}
\end{itemize}
                
\end{frame}
